\documentclass[../../../include/open-logic-section]{subfiles}

\begin{document}

\olfileid{sth}{spine}{Valphabasic}
\olsection{దశల ప్రాథమిక ధర్మాలు}

$V_\alpha$ల నిర్వచనానికి ఉన్న పునాది ప్రాధాన్యాన్ని స్పష్టం చేయడానికి, వాటి గురించి కొన్ని ప్రాథమిక ఫలితాలు ఇస్తాం. ఒక నిర్వచనంతో మొదలెడదాం:\footnote{``potent''కు స్థిరపడిన పదజాలం లేదు; \citet{ButtonLT1} వాడిన పేరు ఇదే.}
\begin{defn}
	$A$ సమితి \emph{సమర్థమైనది} అవుతుంది, అప్పుడూ అప్పుడే $\forall x((\exists y \in A)x \subseteq y \lif x \in A)$.
\end{defn}
\begin{lem}\ollabel{Valphabasicprops}
ప్రతి క్రమసంఖ్య $\alpha$కు:
\begin{enumerate}
	\item\ollabel{Valphatrans} ప్రతి $V_\alpha$ సంక్రమణ సమితి.
	\item\ollabel{Valphapotent} ప్రతి $V_\alpha$ సమర్థమైనది.
	\item\ollabel{Valphacum} $\gamma \in \alpha$ అయితే $V_\gamma
	\in V_\alpha$ (కాబట్టి \olref{Valphatrans} ప్రకారం $V_\gamma \subseteq V_\alpha$ కూడా).
\end{enumerate}
\end{lem}

\begin{proof}
వీటిని (ఏకకాల) అతిపరిమిత ఆగమనం ద్వారా నిరూపిస్తాం. ఆగమనానికి, ప్రతి క్రమసంఖ్య $\beta < \alpha$కు \olref{Valphatrans}--\olref{Valphacum} వర్తిస్తాయని అనుకుందాం.

$\alpha = \emptyset$ సందర్భం సులభం.

$\alpha = \ordsucc{\beta}$ అనుకుందాం. \olref{Valphacum} కోసం, $\gamma \in \alpha$ అయితే ఆగమన ఊహ ప్రకారం $V_\gamma \subseteq V_\beta$; కాబట్టి $V_\gamma \in \Pow{V_\beta} = V_\alpha$. \olref{Valphapotent} కోసం, $A \subseteq B \in V_\alpha$, అంటే $A \subseteq B \subseteq V_\beta$ అనుకుందాం; అప్పుడు $A \subseteq V_\beta$, అందుచేత $A \in V_\alpha$. \olref{Valphatrans} కోసం, $x \in A \in V_\alpha$ అయితే $A \subseteq V_\beta$ కాబట్టి $x \in V_\beta$; ఆగమన ఊహ ప్రకారం $V_\beta$ సంక్రమణ సమితి కావడంతో $x \subseteq V_\beta$, అందువల్ల $x \in V_\alpha$.

$\alpha$ సీమా క్రమసంఖ్య అనుకుందాం. \olref{Valphacum} కోసం, $\gamma \in \alpha$ అయితే $\gamma \in \ordsucc{\gamma} \in \alpha$; ఊహ ప్రకారం $V_\gamma \in V_{\ordsucc{\gamma}}$, కాబట్టి $V_\gamma \in \bigcup_{\beta \in \alpha} V_\beta = V_\alpha$. \olref{Valphatrans}, \olref{Valphapotent} కోసం, సంక్రమణ (వరుసగా సమర్థ) సమితుల సమ్మేళనం కూడా సంక్రమణ (వరుసగా సమర్థ) సమితి అని గమనిస్తే చాలు.
\end{proof}

\begin{lem}\ollabel{Valphanotref}
ప్రతి క్రమసంఖ్య $\alpha$కు $V_\alpha \notin V_\alpha$.
\end{lem}

\begin{proof}
అతిపరిమిత ఆగమనం ద్వారా. స్పష్టంగా $V_\emptyset \notin V_\emptyset$.

$V_{\ordsucc{\alpha}} \in V_{\ordsucc{\alpha}} = \Pow{V_\alpha}$ అయితే $V_{\ordsucc{\alpha}} \subseteq V_\alpha$; ఇంకా \olref{Valphabasicprops} ప్రకారం $V_\alpha \in V_{\ordsucc{\alpha}}$ కనుక $V_\alpha \in V_\alpha$. దీని వ్యతిరేక రూపంలో, $V_\alpha \notin V_\alpha$ అయితే $V_{\ordsucc{\alpha}} \notin V_{\ordsucc{\alpha}}$.

$\alpha$ సీమా క్రమసంఖ్య కాగా $V_\alpha \in V_\alpha = \bigcup_{\beta \in
\alpha}V_\beta$ అయితే, ఏదో $\beta \in \alpha$కు $V_\alpha \in V_\beta$; అయితే $V_\beta \in V_\alpha$ కూడా, అందువల్ల \olref{Valphabasicprops}ను రెండుసార్లు వాడితే $V_\beta \in V_\beta$. దీని వ్యతిరేక రూపంలో, ప్రతి $\beta \in \alpha$కు $V_\beta \notin V_\beta$ అయితే $V_\alpha \notin V_\alpha$.
\end{proof}

\begin{cor}
ఏ క్రమసంఖ్యలు $\alpha, \beta$కైనా: $\alpha \in \beta$ అప్పుడూ అప్పుడే $V_\alpha \in V_\beta$.
\end{cor}

\begin{proof}
ఒక దిశను \Olref{Valphabasicprops} ఇస్తుంది. తిరిగి, $V_\alpha \in V_\beta$ అనుకుందాం. \olref{Valphanotref} ప్రకారం $\alpha \neq \beta$; అలాగే $\beta \notin \alpha$, ఎందుకంటే అలా అయితే $V_\beta \in V_\alpha$; అప్పుడు \olref{Valphabasicprops}ను రెండుసార్లు వాడితే $V_\beta \in V_\beta$, ఇది \olref{Valphanotref}కు విరుద్ధం. కాబట్టి త్రిభేద ధర్మం ప్రకారం $\alpha \in \beta$.
\end{proof}
\noindent
ఇవన్నీ ప్రతి $V_\alpha$ను స్థాయిక్రమంలోని $\alpha$వ దశగా భావించడానికి అనుమతిస్తాయి. ఎందుకంటే:

క్రమసంఖ్యల వాడకం పునరావర్తన భావనను అనుసరిస్తుంది కాబట్టి, మన $V_\alpha$లు ఒక \emph{పునరావర్తన} ప్రక్రియలో ఏర్పడినవిగా భావించవచ్చు. అంతేకాక, ఒక దశ మరొకదానికి ముందు ఏర్పడితే, అంటే $V_\beta \in V_\alpha$, అంటే $\beta \in \alpha$ అయితే, $V_\beta \subseteq V_\alpha$ కాబట్టి నిర్మాణ ప్రక్రియ \emph{సంచితమైనది}. చివరగా, ఏ ఉత్తరవర్తి దశ $V_{\ordsucc{\alpha}}$ అయినా తన ముందు దశ $V_\alpha$ యొక్క ఘాత సమితి కాబట్టి, ముందరి దశలలో అందుబాటులో ఉన్న సమితుల \emph{అన్ని} సాధ్యమైన సమాహారాలను నిజంగానే ఏర్పరుస్తున్నాం.

సంక్షిప్తంగా: $\ZFminus$తో, సమితుల సంచిత-పునరావర్తన స్థాయిక్రమపు మన చిత్రణ \emph{దాదాపు} పూర్తయింది. (అయితే ప్రతిస్థాపనను ఇంకా సమర్థించాలి.)

\end{document}
